% basic nastaveni
\documentclass[12pt]{article}
\setlength{\parindent}{0pt}
\setlength{\parskip}{1em}
\usepackage[utf8]{inputenc}
\usepackage[czech]{babel}
\usepackage[T1]{fontenc}
\usepackage{graphics}
\usepackage{caption}
\usepackage{hyperref}
\usepackage{xcolor}
\usepackage{float}

% matematicke package
\usepackage{amsmath}
\usepackage{amsfonts}
\usepackage{amssymb}
\usepackage{geometry}
\geometry{margin=2.5cm}
\usepackage{pgfplots}
\pgfplotsset{compat=1.18}
\usepackage{mathrsfs}

\title{Maturitní otázka číslo 18: Analytická geometrie lineárních útvarů}
\author{}
\date{}

\begin{document}
\maketitle

\subsection*{Geometrie v rovině}
\subsubsection*{Parametrické vyjádření přímky}
Přímka se dá určit pomocí dvou bodů, které oba protíná a nebo pomocí bodu s směrového vektoru. To
využívá parametrické vyjádření. Mějme bod $A = [A_x, A_y]$ a vektoru $\vec{u} = (u_x, u_y)$. Každý
bod přímky $X$ se dá vyjádřit jako $X = [A_x, A_y] + t \cdot \vec{u}, t \in \mathbb{R}, \vec{u} \ne
0$. Proměnná t se nazývá parametr. Někdy se také parametrické vyjádření rozepisuje do jednotlivých
souřadnic. Pro $X = [x, y]$ platí $x = A_x + t \cdot u_x$ a $y = A_y + t \cdot u_y$.

Pokud bychom chtěli vyjádřit polopřímku, paramet t by byl v rozmezí od $t \in \mathbb{R}, \langle 0,
\infty)$ nebo $(- \infty, 0 \rangle$. Pro vyjádření úsečky by byl parametr $t \in \mathbb{R}, \langle
0, 1 \rangle$. kde by bylo bod $X = [A_x, A_y] + \vec{u}$ druhým bodem úsečky.

\subsubsection*{Obecná rovnice přímky}

\end{document}
